\section{Paths}\label{sec:11-path-algebras:02-paths:1}

\subsection{Recall}\label{sec:11-path-algebras:02-paths:0b}

Formal definition cited in this section, gathered here for quick reference (full citation in the \hyperref[sec:root:bibliography:1]{Bibliography}):

\begin{itemize}
\tightlist
\item
  \textbf{Path.} ``A path of length \(l \geq 1\) with source \(a\) and target \(b\) ... is a sequence \((a \mid \alpha_1, \alpha_2, \ldots, \alpha_l \mid b)\) ... We also agree to associate with each point \(a \in Q_0\) a path of length \(l = 0\), called the trivial or stationary path at \(a\), and denoted by \(\varepsilon_a = (a \| a)\)'' (\cite{AssemSimsonSkowronski2006}, Ch. II §1, p.~42--43, unnumbered definition immediately preceding Definition 1.2 --- which itself defines the \emph{path algebra} built from these paths, not the path concept).
\end{itemize}

A \textbf{path} in \(Q\) is a sequence of arrows that can be composed head-to-tail: \(\alpha_1, \alpha_2, \dots, \alpha_n\) with \(t(\alpha_i) = s(\alpha_{i+1})\) for each consecutive pair. In addition, for each vertex \(i\), a \textbf{trivial path} \(e_i\) of length \(0\) is allowed, starting and ending at \(i\) and composing with nothing but itself.

In the running example, the paths are: \(e_1, e_2, e_3\) (trivial), \(\alpha\), \(\beta\), and \(\beta\alpha\) (first \(\alpha\), then \(\beta\)), and nothing else, since there is no arrow out of \(3\) or into \(1\).

\subsection{References}\label{sec:11-path-algebras:02-paths:2}

Full citations in the \hyperref[sec:root:bibliography:1]{Bibliography}. Formal definitions are gathered in Recall, above.

\begin{itemize}
\tightlist
\item
  Assem, Simson, and Skowroński (\cite{AssemSimsonSkowronski2006}), Ch. II §1, pp.~42--43, unnumbered definition preceding Definition 1.2 --- \textbf{Correction:} this book previously mislabeled the path definition itself as ``Definition 1.2''; that numbered definition is actually the \emph{path algebra} \(KQ\) built from these paths, not the path concept, which is unnumbered text immediately before it.
\item
  Schiffler (\cite{Schiffler2014}), \textbf{Definition 2.1 and Example 2.2} (Chapter 2, §2.1) --- same notion, called the ``constant path'' (or ``lazy path'') \(e_i\) at vertex \(i\) for the length-\(0\) case.
\end{itemize}
